\documentclass[11pt]{amsart}
\usepackage[T1]{fontenc}
\usepackage{newtxtext,newtxmath}
\usepackage[margin=1in]{geometry}
\usepackage{booktabs,fancyvrb}
\usepackage[hidelinks]{hyperref}
\raggedbottom
\newtheorem{theorem}{Theorem}[section]
\newtheorem{lemma}[theorem]{Lemma}
\title[Two certified prime gaps]{Consecutive primes at distances 123,662 and 125,318}
\author{Will Howe}
\email{howew@purdue.edu}
\author{Grant Huffman}
\email{huffma97@purdue.edu}
\author{Jensen Kohlmeyer}
\email{jenwin@purdue.edu}
\author{Liam Kruer}
\email{kruerl@purdue.edu}
\date{November 16, 2026}
\subjclass[2020]{Primary 11A41; Secondary 11Y11, 11Y55}
\keywords{Prime gaps, covering congruences, primality certificates}
\begin{document}
\begin{abstract}
We exhibit consecutive-prime gaps of exactly $123\,662$ and $125\,318$,
with endpoints of $8\,416$ and $6\,454$ decimal digits, respectively.
Covering congruences certify most interior integers as composite. For the
second gap, the remaining integers are certified by proper divisors or
Fermat witnesses. Elliptic-curve certificates prove all four endpoints prime.
We give the elementary arguments that turn these finite arithmetic checks
into a proof. The proof is computer-assisted.
\end{abstract}
\maketitle

\section{The two gaps}
For $h\in\{123\,662,125\,318\}$, let $p_h$ be the integer displayed in
Appendix~A or~B, respectively.
\begin{theorem}\label{thm:gaps}
For each of these two values of $h$, the integers $p_h$ and $p_h+h$ are
consecutive primes. Both endpoints have $8\,416$ digits when $h=123\,662$
and $6\,454$ digits when $h=125\,318$.
\end{theorem}
The full covering tables, additional compositeness witnesses, and four
primality chains form the supplementary arithmetic appendix. Their mathematical
format and the required checks are described below.

For each gap, the covering table consists of pairs $(q,b_q)$ with distinct
prime moduli $q$ and $b_q\not\equiv0,-h\pmod q$. Define
\[
 W_h=\prod q,\qquad 0\le a_h<W_h,\qquad a_h\equiv b_q\pmod q.
\]
The Chinese remainder theorem determines $a_h$ uniquely, and the displayed
starting integers satisfy $p_h=a_h+m_hW_h$, with the following parameters.
\begin{center}
\begin{tabular}{r r r r r}
\toprule
$h$ & $m_h$ & Number of moduli & Largest modulus & Digits of $W_h$\\
\midrule
$123\,662$ & $1\,523\,845$ & $2\,216$ & $32\,587$ & $8\,410$\\
$125\,318$ & $5\,046\,505$ & $1\,754$ & $14\,983$ & $6\,448$\\
\bottomrule
\end{tabular}
\end{center}
These progressions were obtained by optimizing the covering congruences.
Sieving and endpoint tests then found the displayed multipliers.

\section{Compositeness of the interior}
If $1\le j<h$ and $j\equiv-b_q\pmod q$, then $q\mid p_h+j$.
Since $1<q<p_h+j$, this is a proper divisor. The exact covering counts are
\begin{center}
\begin{tabular}{r r r r}
\toprule
$h$ & Covered by congruences & Additional divisors & Fermat witnesses\\
\midrule
$123\,662$ & $123\,661$ & $0$ & $0$\\
$125\,318$ & $124\,526$ & $175$ & $616$\\
\bottomrule
\end{tabular}
\end{center}
For the second gap, the remaining $791$ offsets are listed in the arithmetic
appendix. At $175$ of them, an integer $d$ satisfies $1<d<p_h+j$ and
$d\mid p_h+j$. At each of the other $616$, direct modular exponentiation gives
\[
 2^{p_h+j-1}\not\equiv1\pmod{p_h+j}.
\]
Fermat's little theorem shows that these integers are composite. The three
columns account for every interior offset, proving the required
compositeness in both cases.

\section{Primality of the endpoints}
We use the elliptic-curve criterion underlying ECPP~\cite{atkinmorain}.
\begin{lemma}\label{lem:ecpp}
Let $N>3$ be coprime to $6$ and let $E:y^2=x^3+Ax+B$ satisfy
$\gcd(4A^3+27B^2,N)=1$. Suppose an affine point $R$ satisfies $[r]R=O$,
where $r$ is prime and $r>(N^{1/4}+1)^2$. If every division in the
computation has denominator coprime to $N$, then $N$ is prime.
\end{lemma}
\begin{proof}
If $N$ were composite, it would have a prime divisor $\ell\le\sqrt N$.
Reduction modulo $\ell$ gives a nonsingular elliptic curve and a nonzero
point of order $r$. Hasse's bound gives
$r\le\#E(\mathbb F_\ell)\le(\sqrt\ell+1)^2\le(N^{1/4}+1)^2$,
a contradiction.
\end{proof}

Each endpoint certificate is a chain of tuples $(N_i,t_i,s_i,A_i,X_i,Y_i)$.
Set $N_{i+1}=(N_i+1-t_i)/s_i$, $P_i=(X_i,Y_i)$, and
$B_i\equiv Y_i^2-X_i^3-A_iX_i\pmod{N_i}$. At every step the exact checks are
\begin{equation}\label{eq:checks}
\begin{gathered}
 N_i>3,\quad s_i>0,\quad N_{i+1}\in\mathbb Z,\quad
 \gcd(N_i,6)=\gcd(4A_i^3+27B_i^2,N_i)=1,\\
 N_{i+1}>(N_i^{1/4}+1)^2,\qquad
 R_i=[s_i]P_i\ne O,\qquad [N_{i+1}]R_i=O.
\end{gathered}
\end{equation}
The points $R_i$ are affine, and all divisions in these group operations
have denominators coprime to $N_i$. The four chains end as follows.
\begin{center}
\begin{tabular}{l r r}
\toprule
Starting integer & Steps & Terminal prime\\
\midrule
$p_{123\,662}$ & $679$ & $7\,440\,766\,535\,021$\\
$p_{123\,662}+123\,662$ & $671$ & $636\,419\,037\,730\,097$\\
$p_{125\,318}$ & $538$ & $450\,463\,843\,321$\\
$p_{125\,318}+125\,318$ & $537$ & $403\,307\,424\,311\,829\,421$\\
\bottomrule
\end{tabular}
\end{center}
For completeness, the arithmetic appendix gives elementary order
certificates for the four terminal primes. Each row specifies $N$, a
complete factorization of $N-1$, and an integer $g$ such that
\[
 g^{N-1}\equiv1\pmod N,\qquad
 \gcd\bigl(g^{(N-1)/q}-1,N\bigr)=1\quad(q\text{ prime},\ q\mid N-1).
\]
For every prime divisor $\ell$ of $N$, these checks force $g$ to have order
$N-1$ modulo $\ell$. Thus $N-1\mid\ell-1$, so $\ell\ge N$ and $N$ is prime.
The factor primes are certified recursively in increasing order, starting
with primes below $100$, which are checked by trial division.

\begin{proof}[Proof of Theorem~\ref{thm:gaps}]
Starting from the terminal primes and applying Lemma~\ref{lem:ecpp}
backwards along each chain proves all four endpoints prime. Section~2
proves every integer between the corresponding endpoints composite.
The displayed decimal expansions and addition of $h$ give the stated
digit counts.
\end{proof}

\subsection*{Acknowledgment}
The authors acknowledge OpenAI Codex for assistance with the construction,
finite arithmetic verification, and preparation of this paper.
\begin{thebibliography}{1}
\bibitem{atkinmorain}
A.~O.~L. Atkin and F. Morain,
\emph{Elliptic curves and primality proving},
Math. Comp. \textbf{61} (1993), no.~203, 29--68.
\end{thebibliography}

\appendix
\clearpage
\section{The starting prime for $h=123\,662$}\label{app:first}
The following lines concatenate to $p_{123\,662}$. Spaces separate blocks
of ten digits. The right endpoint is obtained by adding $123\,662$.
\begin{Verbatim}[fontsize=\fontsize{6.4}{7.3}\selectfont]
1978533019 1033653481 5536674032 5124196129 8851048675 9715711486 7089520013 2794096239 6866825311 1422429041 3682734817 7817220447
5014214217 4536286439 8570782519 1235789485 7383945998 1528781234 4848619129 5146200069 6664140930 6935569934 3821784681 2770631783
5359989451 9959718327 4683348047 2188926152 6363302453 2274977941 8302494490 9769861992 6187829493 1931840291 7753411792 2888025483
7385423067 3816189600 3475050237 9074870230 5873592406 1122624139 0687369215 9609949041 3642675456 8960183613 4506953214 0133457686
3258896767 2295270948 4528077222 2028775473 6086047244 0609189562 4627341893 6160871890 1467859905 0063814165 3358572956 8693394831
0699055811 8785872261 5964945103 2023673658 5624385095 9287886616 9175610393 5440274564 6129550704 1387226846 3541788609 6661626751
4060433495 0442415942 9149404097 6238812364 7274302786 4653744906 5948899455 1588980006 4217426901 7443088639 2298121644 5902682980
6487663658 8951049405 6071898675 8617428664 5002524377 6981349802 1930628070 2244184502 9672486600 1669349989 2927951238 8907227496
9687383412 5633161062 9226422550 7979506704 3741160088 6048541842 9204870425 0020410417 8518688592 5289370412 6289671903 5401547238
8764999563 3571649027 7090718455 4753708618 5242849441 8533137181 1845240610 3362560113 1686238117 5114492873 1934514832 1082597006
5223557572 3941734115 5673864377 9749730348 0479459509 5755066676 5515154518 1174112778 7767326587 5575737683 9492662185 1477273873
7349145971 2859690431 2573886037 6522313803 2485332786 2545847368 0748787503 6851066014 4114598998 8137848021 7874057180 8004094395
0018998937 2931107593 3977416116 4243234619 8293362555 4849916138 7612169975 4798121359 1063824136 2158181430 8813776086 9183720845
7665991130 5601828600 5242629352 6760664907 1812934160 7361318741 1669785965 8375512724 7777713745 1528880348 2592427504 9226793025
5245074256 0675278179 9333627233 5961663108 4230149770 0760935689 4844331722 0665941655 5673930843 6291099542 3065235457 1776421220
2795221902 8052926700 8567044528 6152296339 3306839914 1297415021 1196980356 0243186835 1548739124 7559271194 1783876788 0448422375
9198976489 9028143434 4209510869 7488957359 2355367615 2601091202 2197802600 2983844446 9722825191 9668926737 9114121005 5630216254
5192204036 7470381949 1719567093 0481348956 9909777979 7966031767 1298071029 8193551446 7597114922 4537324987 0831629846 3931563057
8247162489 9648874089 8273876652 2440692322 2085975035 4285888987 6748311726 0243865865 1118024765 3380652506 2163946458 9273665798
6016535913 4837350950 2415453259 0961270445 9428286060 5957698558 8602682341 2823486259 0641922700 3731433630 3818888190 9840092349
0857369460 1479921253 3684963056 6067301711 0859712237 7693690577 5694798622 1409835016 9765352099 9562454017 8924739665 7000534704
1424048758 5375857278 8789805899 4312117184 9756700530 5454900558 4330578303 4624323837 5050158639 9037777428 7996485602 5752341511
1321131749 0929451390 5129002316 6263750984 8426820090 2593752716 5963747079 0109699686 7068709080 2235663050 6371097773 4658448604
7081492012 6345770352 9471022457 1224664611 3033530882 2896907816 4809735026 9065990144 2877959491 5880628381 5869810694 6302824357
0421270443 6400774047 2265697248 5559773082 8111879454 3715106854 6476476890 2531672319 9679346080 1999125110 3788595248 6396902876
6785168423 2260975808 9623352725 2768200030 2927505119 0142322699 2913643117 1150909265 3986526327 3331277189 8185531743 7415355083
2924570870 2144738705 4266794124 4797675249 6077811687 1781956956 9252609689 1546443177 9878024291 9515377354 4725906265 9824964975
3426221641 4787460234 9459976194 6705092906 7025479043 3388967398 0439657058 1041556179 3620743348 2615080736 1049073589 4622985662
6871509156 8643519676 0997543582 9037980241 4232284925 2708820173 3202806501 3687113498 0678966594 9391280000 9576330909 4075062165
0277087670 9758938136 9042552710 3877034458 5707105565 9144798597 6853527099 3527232939 8533072969 2887754577 1248202938 6954105840
6158224081 9595285625 9968040851 5368748211 6424410633 7215184636 1235631607 0734416684 5591213922 2043107940 2016219039 7832300379
7523047885 2760431063 4699586806 6138126873 0350706054 4134921073 3532795531 2753552133 0201630174 4568949902 5070480000 4206543281
2917830680 2739268272 5908146871 9476548741 3101685817 2626159488 7735531466 1747335457 9056445165 7498922148 1544604832 8364588091
3639246728 5927470803 8640798494 6491520092 2629457379 5110762460 8594410225 4488625856 1833882148 2746393139 2176804107 6266704171
7712772644 7980541068 0735526955 8157242543 7458727553 6602037116 8605250490 2255328206 5967488338 5060236414 7633814049 6619293006
8802763554 0342016269 2450495306 7126979619 6376835294 1035332701 9152818633 7212722439 4448400142 1517353543 7572055462 7408795597
2352960979 9767406970 4464978064 3667040571 1941567220 2855945032 5753296580 5677411139 9014564031 2477045944 9850885348 8996335848
9304637603 0376800804 1828474128 2784465485 8621439135 2307580293 2612685343 7471548500 2275795291 9075154648 0634992641 7261296756
2257977637 7223589629 2089505569 1616615033 9861554507 4626013537 7306181987 5244805823 3373736630 0295648593 3317984682 2238871942
8439310152 2524951635 8768731544 7436352405 1994378408 7754984133 2436768574 1387839312 5892642527 6855527538 0523066523 6181497644
7346925768 8264749060 4695404101 0269691631 8375658221 1845085582 4845081050 1631945578 1004148111 5713172862 6075774963 2487692626
5951402006 3149203290 5162164977 2743142501 8939728245 4086097889 6230500094 3311706633 1268635971 0714088975 1434911181 0366367692
2283023442 2440478259 1775457479 5766960600 5968956759 3403212758 7934879603 8042978769 4074440159 6583038078 9643381856 4977315225
4647626359 1383249630 4892691098 7048814494 4105111287 6722879777 0865767689 6321581863 1709182205 7650566498 7631662459 2830220628
2364792166 5058824508 9251811324 5580621218 7807859254 2996898941 1657994688 4629380118 4014145467 5855274309 3603569329 2032330510
5977300417 2488915682 5546625108 4053553570 4031169746 5294442327 0773334264 3040457497 6139860630 8120318163 1217964451 7989663018
9092019834 6951338258 1396327334 0704716994 8279538679 1511312408 5370780214 2853619811 7867493487 8321111872 1452514559 7313228717
0006538536 4676766584 1910675337 0413746229 0646628689 9122345144 4881247382 6551606895 2118544715 1324702342 6960334957 6406031628
5369721568 3910294416 2524692816 4607062780 5757014917 0060617272 7668526903 5100452069 4050853864 1457993769 7857809209 6453914181
5010676540 3304273321 5639372054 0322413992 7603961715 6596014640 2675995628 5937287878 7661976441 4854932940 1485718041 0116256597
7481246703 1511212391 1060564170 8459010478 2964857233 7202799436 9444988146 0022889750 3102908736 8162917943 8201381579 1661602919
3522500892 7935304437 9612892715 8973102531 8295933662 0637687031 3912461793 1660865184 6016590503 8165409311 4400721005 2131804156
4751835172 8033114345 4826818994 8274471188 8281617761 7309142596 1229433393 8762914626 7848667754 5258006541 6457532134 8086573750
2177921607 7573455720 0121744042 2256115427 2280454618 1903913927 2938375112 2581571556 8189105348 9698480559 0790811359 0020132851
4439880534 8414144195 5404114573 1498483252 2701461141 4034619034 1743645179 2753698944 7926752008 4735057624 3192489821 3407011321
1056285873 7416814133 9689470402 6511940284 0758477710 8365202543 5959871717 4009919572 1681012234 8811001964 6289063303 9670710201
0286788768 8141647910 0423460557 9896540285 8178970002 3365180913 0257431263 9402847815 3614480745 1994120842 5708557357 3689754099
7347510145 8602475308 8242361395 0279229633 0962119790 4689949160 0164467521 9785288191 2589964675 4394443621 6173801356 1261385794
4059378488 3890675067 2606322658 8255133419 8739752266 3791911279 0224061274 7776136683 8765764279 3910808863 5316171658 7695862115
3248829866 1012399852 7605459960 2943768132 9000043016 9914840491 3467556069 5626988951 9470706692 6410090672 0139576985 1276936028
6277614613 3908310209 4752344046 3101460371 1819830523 1298109280 1906452116 9445262148 7020580934 3323658121 2392370371 0064795759
8213655263 9121742817 2718501775 1349331363 6139942029 5528707066 7581861477 3338533016 1801800394 7347517498 8262701210 9907086824
1565655241 6237567758 2995819637 2361299306 7775821900 1500235171 6821601874 4444631733 6406109588 4387927328 4750663536 4886136437
3686548179 4281785305 4782681604 4381753423 6972860913 1673626680 8773935373 4883515901 3376827313 4850406055 3070614174 3525867465
8751024118 1960531396 9666522677 8381272702 7452228100 7809397227 3322181561 7753567266 4341225364 9464961317 2842795095 1609195218
6665265024 0510063303 7595982122 5027411783 3398384598 0963913202 6587723023 8206878362 2471078074 5579831967 5984833098 4762370655
8396450640 3650199517 6695282997 7336578026 7556001238 0623354164 8249450542 0939107243 7818065784 5771760286 7209113747 5281535627
5367577256 7460186654 5900311339 6796892563 5156021599 1748794840 0571543572 6140760930 5200775590 2987066249 0459564822 4555782326
9463395032 4077634960 7982271336 6862296857 8509405265 5721471232 4539720707 3674509316 4392738043 5621561245 7847240840 8625161628
9865300046 1294199498 4026084184 6749988697 3762950016 3082003277 6153838005 1679679976 8092530840 6230099904 9643743672 5611218287
2815090479 351597
\end{Verbatim}

\clearpage
\section{The starting prime for $h=125\,318$}\label{app:second}
The following lines concatenate to $p_{125\,318}$. Spaces separate blocks
of ten digits. The right endpoint is obtained by adding $125\,318$.
\begin{Verbatim}[fontsize=\fontsize{6.4}{7.3}\selectfont]
5177351185 4459854632 7695956438 5729587744 1335282968 5638645816 7233145552 3593544663 3379922985 5373108128 0426918899 1827237135
4007000473 7048513604 1957567607 5628001399 9552480709 5030662041 0684998996 4759223834 0809705674 8593389081 4689232124 1493364031
2369256582 3051678914 8017442137 9665363916 2238276540 5267470761 6276757828 7614073744 1375150193 6024938797 2062830705 0685798950
8800353734 8756451343 9598753295 9428991260 0800705654 7490321289 5472619402 0451583772 3252614117 4492155074 8341990009 1849481618
9609864853 1881085580 2271933357 8495936666 0722748165 3258682746 5905588655 9438031737 8499704134 2378077781 0563679596 3346938323
6805297659 7722615170 1333031409 4252845874 7550553495 4026250327 9870560003 5390155917 4397922503 8028846988 3005585380 8606665496
3940320482 9684065357 7589265064 2381804726 3158185121 7617216056 6550113974 6001018405 4295011674 8258709519 7274583649 4908002644
9841160374 6875315234 9917858330 0714937242 4411139276 4424469029 8687147673 3698136225 2021143382 5556041820 9276526538 7697243942
6678789614 2383456323 0198235411 1274212725 5585426191 6193743735 7888797105 8946446609 3972929417 6791627905 0047425613 0274317818
7112385335 1161902171 2086800694 1643903182 2966218142 6953879713 8971622201 7470670207 5648945669 7389974790 7837222538 6005173422
6835735565 6301173647 0320909714 9326968134 4325641064 6308115045 8107577472 9348130816 7668863905 2862717712 6636521766 3689457354
7888481549 5585973695 2551903844 5377571347 6072556333 6600701191 6032218848 2315516632 0677001801 1270312331 6475698800 2724343249
3484515265 7453837988 5061412971 7187758279 7663369055 1132434579 2980458350 2256998838 1406827288 8941016926 6863590673 5868843678
9619493363 4560775913 9982060145 9312735662 1210094811 6622881224 5885658483 8826850741 4346815783 2123515263 1309233520 9794185086
9226268525 5399480296 5090395730 3516366326 8556534773 6178880303 5449635406 7910621380 9354252778 2160790690 9006754671 0120937392
1654230275 8617466504 3868300800 7411370859 8689189418 2113362216 1499746736 1214518621 3758413483 2781492568 1144698455 7257818142
4971091147 8355740306 4046441833 1973039070 3234576913 3191911765 8511312184 1297363553 4964506426 6737408424 0475484864 9816893005
6728560503 9363007766 2231075597 3566590394 0381863555 2549422078 6969555064 9340537613 2044964599 6545313680 3379309107 4946744178
6301365383 2898432389 5409152801 1039312004 6858019336 0096019125 1264980634 1039537093 9016180728 7070581519 5137520132 0153133318
8934856608 6997206774 8200079388 6971906837 4539137143 2444232263 9700450492 9440326313 6828288040 1427603096 7115773974 5970240443
2461605149 6135212260 3704185649 1840492059 7756182571 6668015147 6424380881 7766799933 0018225904 6414668288 9007761386 5826024267
8836958515 8670957310 1816120695 7580699692 5240000463 2363035244 7871694808 4783735603 2492041407 3908374693 8198088775 4356008766
1749837106 5495794170 3686637313 0107614530 4933980448 1975084923 6041706382 1121522515 0153664766 8569357553 5072459756 9117529849
6479543363 8110921369 5905568221 5547356978 1349214041 4385606285 6259653328 7043533230 2586147463 0548240641 1794845737 5979418185
1297001764 8272126550 7474758177 1014796710 9902694088 8879879429 0568526930 6634168647 4343141070 1043208850 3444300107 3494180786
1324716684 7188787817 3299761455 4340221994 0429348301 3001262241 3130406607 0872155862 1233785846 9631554864 7360607019 6029450783
3183548958 5160936355 9948330428 6395166663 5756384200 2553772901 1256111154 7841166898 0112581437 3329431022 1900195404 8826240986
9103105231 4397051501 6818999533 6880157105 5944253023 4063928086 7168156995 4506596472 4077542189 7095009404 6651831313 9446296494
5175470030 3227765706 3247927992 8084105191 7986271145 5561975903 9726060147 4754501581 9282120737 6297642387 9089037711 3043625464
4534763541 0641205654 8894214364 4352365768 2758755096 0036469531 7994778857 3590416659 8474229224 0464459313 2066795569 0309155475
0237609297 0230581594 4316986555 8231744734 8905416731 6154243102 1251746175 7247964698 6281656112 0340205775 6177242498 0856513903
2762573552 8929687971 2822598939 3704764985 6637377461 2742908550 0008095627 3712897333 3324626490 5237715665 2194249239 9249140394
8877146619 7410856997 5434093467 0501380475 5733907816 1034017204 9827896417 5602972343 6563395448 2810095079 6329953291 8179387410
3247414242 7925528000 7559712718 1701103312 4869096557 5281723203 4331656158 3981854131 4219642015 0757711473 3025981957 0992677489
8934862729 4800722582 5577064350 9538424200 1090799929 8001207673 5710257936 4088517256 4534251977 4000192363 7436370580 8732459658
1540928214 6368880146 3480217656 0619362792 9619294914 1515413593 0801233578 3287195757 3356326584 1014759181 4122701683 3123745314
1002703806 4516542597 9010342108 5098363045 0584049699 9908775883 5781091903 7833862838 9411886979 2814364834 0067352011 2446086701
4020443282 1003207622 3715151890 0721603760 4707757875 3651803604 6234442089 8652989140 1607123789 4790018990 1814850737 1823528745
4184504004 8448663553 0624842415 2303541547 3960660024 9773515910 1096479673 4769562957 3149546039 3910817630 8559270586 5147852077
1125121576 7719136621 3440668762 9992111957 6038350009 2028158498 0867754892 5818518518 7544588239 4563919511 8563053916 7290225274
3458520931 5061729937 7381376161 6622912510 0114385946 9802204149 9470037274 9299497845 8988324657 6485975344 2917616595 6609233835
0885924873 4969939368 3318330522 8859740426 3007439101 7603875163 7581045027 9737396740 9694522227 5841102785 4106047319 3559714947
7616908508 0994647581 3007865665 5470828797 4922162163 9108887522 7378750514 8075947102 0553713332 9456034536 3799375584 2283355929
0685690104 8459597706 5140567351 8650942920 1425875196 4480766461 4818188189 0351891455 4925017808 9325218474 2072566505 6636394810
0405031764 7717598293 5040910817 5099821434 9524951698 6682606645 5090880720 0297151816 2354135843 9335965418 4541674332 9934642250
8208346340 6125091488 0733317654 1146675280 2266836807 5639312465 8150630815 9563727958 4679161250 5484322134 4506166334 7518642690
5238083860 3327939261 4832546833 1502253898 4814095257 8233935146 0952094064 1432973832 9337735075 2014941985 7466877493 6824578000
2199326759 1397807700 9259937928 9462918821 2641677697 5703562685 5059807344 8312803382 0378507082 1837553458 7156490929 0497303292
0709280323 5525043000 1136726874 3645511943 0529677769 3290347498 6307799338 5740332956 2077598921 3910869077 4392440681 1446578525
0411582359 7485142692 9139034601 6799445646 7772302625 0983180020 3428796593 6032119930 0187281074 1920971837 5746064136 5330242499
0336025569 3233140697 8038463225 1266408180 1304813660 6047769457 7810311344 6950202370 9519023896 8495945079 9315570869 5617217798
5121978690 2985326947 6381644819 9185385979 4268325492 1219704311 4040258985 5347163177 3112317737 5034618296 3214398142 9729425375
5882899223 8246724138 8357029791 6486893309 2315213708 7166536322 0535280189 1899146492 0085098042 1592514763 0767703164 6157512995
8699106062 4663612956 8166669842 8323014535 7919493586 5593787725 3967041963 0812470815 6805772741 8089
\end{Verbatim}

\end{document}
